\SourceLeaf{61}%
% A544-P0061-B02
\begin{center}
\textit{THEOREMATIS FUNDAMENTALIS IN THEORIA RESIDUORUM QUADRATICORUM}\par
\textit{DEMONSTRATIO QUINTA.}
\end{center}

% A544-P0061-B03
\Article{1}

% A544-P0061-B04
In introductione iam declaravimus, demonstrationem quintam et tertiam ab eodem lemmate proficisci, quod commoditatis caussa, in signis disquisitioni praesenti adaptatis hoc loco repetere visum est.

% A544-P0061-B05
\textsc{Lemma.}\quad \textit{Sit $m$ numerus primus (positivus impar), $M$ integer per $m$ non divi}\textit{sibilis; capiantur residua minima positiva numerorum}

% A544-P0061-B06
\[
\begin{aligned}
M,\ 2M,\ 3M,\ 4M\ \mathinner{.\,.\,.\,.\,.\,.}\ \tfrac12(m-1)M
\end{aligned}
\]

% A544-P0061-B07
\textit{secundum modulum $m$ quae partim erunt minora quam $\tfrac12m$, partim maiora: poste}\textit{riorum multitudo sit $=n$. Tunc erit $M$ residuum quadraticum ipsius $m$, vel non}\textit{residuum, prout $n$ par est, vel impar.}

% A544-P0061-B08
\textsc{Demonstr.}\quad Sint e residuis illis ea, quae minora sunt quam $\tfrac12m$, haec $a,b,c,d$ etc., reliqua vero, maiora quam $\tfrac12m$, haec $a',\ b',\ c',\ d'$ etc. Posteriorum complementa ad $m$, puta $m-a',\ m-b',\ m-c',\ m-d'$ etc. manifesto cuncta minora erunt quam $\tfrac12m$, atque tum inter se tum a residuis $a,\ b,\ c,\ d$ etc. diversa, quamobrem cum his simul sumta, ordine quidem mutato, identica erunt cum omnibus numeris $1,\ 2,\ 3,\ 4\dotsfour\tfrac12(m-1)$. Statuendo itaque productum

% A544-P0061-B09
\[
\begin{aligned}
1\pdot2\pdot3\pdot4\dotsfive\tfrac12(m-1)=P
\end{aligned}
\]

% A544-P0061-B10
erit

% A544-P0061-B11
\[
\begin{aligned}
P=abcd\dotsfour\times(m-a')(m-b')(m-c')(m-d')\dotsfour
\end{aligned}
\]

% A544-P0061-B12
adeoque

% A544-P0061-B13
\[
\begin{aligned}
(-1)^nP=abcd\dotsfour\times(a'-m)(b'-m)(c'-m)(d'-m)\dotsfour
\end{aligned}
\]

% A544-P0061-B14
Porro fit, secundum modulum $m$,

% A544-P0061-B15
\[
\begin{aligned}
PM^{\frac12(m-1)}&\equiv abcd\dotstwo\times a'b'c'd'\dotstwo\\&\equiv abcd\dotstwo\times(a'-m)(b'-m)(c'-m)(d'-m)\dotstwo
\end{aligned}
\]

% A544-P0061-B16
adeoque

% A544-P0061-B17
\[
\begin{aligned}
PM^{\frac12(m-1)}\equiv P(-1)^n
\end{aligned}
\]

% A544-P0061-B18
Hinc $M^{\frac12(m-1)}\equiv\pm1$, accepto signo superiori vel inferiori, prout $n$ par est vel impar, unde adiumento theorematis in \textit{Disquisitionibus Arithmeticis} art. 106 demonstrati lemmatis veritas sponte demanat.

\SourceLeaf{62}%
% A544-P0062-C02
\Article{2}

% A544-P0062-C03
\textsc{Theorema.}\quad \textit{Sint $m,M$ integri positivi impares inter se primi, $n$ multitudo} \textit{eorum e residuis minimis positivis numerorum}

% A544-P0062-C04
\[
\begin{aligned}
M,\ 2M,\ 3M\dotsfive\tfrac12(m-1)M
\end{aligned}
\]

% A544-P0062-C05
\textit{secundum modulum $m$, quae sunt maiora quam $\frac12m$; ac perinde $N$ multitudo eorum} \textit{e residuis minimis positivis numerorum}

% A544-P0062-C06
\[
\begin{aligned}
m,\ 2m,\ 3m\dotsfive\tfrac12(M-1)m
\end{aligned}
\]

% A544-P0062-C07
\textit{secundum modulum $M$, quae sunt maiora quam $\frac12M$.\quad Tunc tres numeri $n,N$,} \textit{$\frac14(m-1)(M-1)$ vel omnes simul pares erunt, vel unus par duoque reliqui impares.}

% A544-P0062-C08
\textsc{Demonstr.}\quad Designemus

% A544-P0062-C09
{\small
\begin{center}\setlength{\tabcolsep}{3pt}
\begin{tabular}{@{}lll@{}}
per $f$ & complexum numerorum & $1,\ 2,\ 3\dotsfive\frac12(m-1)\mathord{.}$ \\
per $f'$ & complexum numerorum & $m-1,\ m-2,\ m-3\dotsfive\frac12(m+1)$ \\
per $F$ & complexum numerorum & $1,\ 2,\ 3\dotsfive\frac12(M-1)$ \\
per $F'$ & complexum numerorum & $M-1,\ M-2,\ M-3\dotsfive\frac12(M+1)$
\end{tabular}\end{center}
}

% A544-P0062-C10
Indicabit itaque $n$, quot numeri $Mf$ residua sua minima positiva secundum modulum $m$ habeant in complexu $f'$, et perinde $N$ indicabit, quot numeri $mF$ habeant residua sua minima positiva secundum modulum $M$ in complexu $F'$. Denique designet

% A544-P0062-C11
{\small
\begin{center}\setlength{\tabcolsep}{3pt}
\begin{tabular}{@{}lll@{}}
$\varphi$ & complexum numerorum & $1,\ 2,\ 3\dotsfive\frac12(mM-1)$ \\
$\varphi'$ & complexum numerorum & $mM-1,\ mM-2,\ mM-3\dotsfive\frac12(mM+1)$
\end{tabular}\end{center}
}

% A544-P0062-C12
Quum quilibet integer per $m$ non divisibilis secundum modulum $m$ vel alicui residuo ex $f$ vel alicui ex $f'$ congruus esse debeat, ac perinde quilibet integer per $M$ non divisibilis secundum modulum $M$ congruus sit vel alicui residuo ex $F$ vel alicui ex $F'$: omnes numeri $\varphi$, inter quos manifesto nullus per $m$ et $M$ simul divisibilis occurrit, in octo classes sequenti modo distribui possunt.

% A544-P0062-C13
I.\quad In prima classe erunt numeri secundum modulum $m$ alicui numero ex $f$, secundum modulum $M$ vero alicui numero ex $F$ congrui.\quad Designabimus multitudinem horum numerorum per $\alpha$.

\SourceLeaf{63}%
% A544-P0063-C02
II.\quad Numeri secundum modulos $m,M$ resp. numeris ex $f,F'$ congrui, quorum multitudinem statuemus $=\beta$.

% A544-P0063-C03
III.\quad Numeri secundum modulos $m,M$ resp. numeris ex $f',F$ congrui, quorum multitudinem statuemus $=\gamma$.

% A544-P0063-C04
IV.\quad Numeri secundum modulos $m,M$ resp. numeris ex $f',F'$ congrui, quorum multitudo sit $=\delta$.

% A544-P0063-C05
V.\quad Numeri per $m$ divisibiles, secundum modulum $M$ vero residuis ex $F$ congrui.

% A544-P0063-C06
VI.\quad Numeri per $m$ divisibiles, secundum modulum $M$ vero residuis ex $F'$ congrui.

% A544-P0063-C07
VII.\quad Numeri per $M$ divisibiles, secundum modulum $m$ autem residuis ex $f$ congrui.

% A544-P0063-C08
VIII.\quad Numeri per $M$ divisibiles, secundum modulum $m$ vero residuis ex $f'$ congrui.

% A544-P0063-C09
Manifesto classes V et VI simul sumtae complectentur omnes numeros $mF$, multitudo numerorum in VI contentorum erit $=N$, adeoque multitudo numerorum in V contentorum erit $\frac12(M-1)-N$.\quad Perinde classes VII et VIII simul sumtae continebunt omnes numeros $Mf$, in classe VIII reperientur $n$ numeri, in classe VII autem $\frac12(m-1)-n$.

% A544-P0063-C10
Prorsus simili modo omnes numeri $\varphi'$ in octo classes IX\,.\,.\,XVI distribuentur, in quo negotio si eundem ordinem servamus, facile perspicietur, numeros in classibus

% A544-P0063-C11
\begin{center}IX, X, XI, XII, XIII, XIV, XV, XVI\end{center}

% A544-P0063-C12
contentos resp. esse complementa numerorum in classibus

% A544-P0063-C13
\begin{center}IV, III, II, I, VI, V, VIII, VII\end{center}

% A544-P0063-C14
contentorum ad $mM$, ita ut in classe IX reperiantur $\delta$ numeri; in classe X, $\gamma$ et sic porro.\quad Iam patet, si omnes numeri primae classis associentur cum omnibus numeris classis nonae, haberi omnes numeros infra $mM$, qui secundum modulum $m$ alicui numero ex $f$, secundum modulum $M$ vero alicui numero ex $F$ sunt congrui, quorumque multitudinem aequalem esse multitudini omnium combinationum singulorum $f$ cum singulis $F$, facile perspicitur.\quad Habemus itaque

% A544-P0063-C15
\[
\begin{aligned}
\alpha+\delta=\tfrac14(m-1)(M-1)
\end{aligned}
\]

\SourceLeaf{64}%
% A544-P0064-C02
similique ratione etiam erit

% A544-P0064-C03
\[
\begin{aligned}
\beta+\gamma=\tfrac14(m-1)(M-1)
\end{aligned}
\]

% A544-P0064-C04
Iunctis omnibus numeris classium II, IV, VI, manifesto habebimus omnes numeros infra $\frac12mM$, qui alicui residuo ex $F'$ secundum modulum $M$ congrui sunt.\quad Iidem vero numeri ita quoque exhiberi possunt:

% A544-P0064-C05
\[
\begin{aligned}
F',\ M+F',\ 2M+F',\ 3M+F'\dotsfour\tfrac12(m-3)M+F'
\end{aligned}
\]

% A544-P0064-C06
unde omnium multitudo erit $=\frac14(m-1)(M-1)$, sive habebimus

% A544-P0064-C07
\[
\begin{aligned}
\beta+\delta+N=\tfrac14(m-1)(M-1)
\end{aligned}
\]

% A544-P0064-C08
Perinde e iunctione omnium classium III, IV, VIII colligere licet

% A544-P0064-C09
\[
\begin{aligned}
\gamma+\delta+n=\tfrac14(m-1)(M-1)
\end{aligned}
\]

% A544-P0064-C10
Ex his quatuor aequationibus oriuntur sequentes:

% A544-P0064-C11
\[
\begin{aligned}
2\alpha&=\tfrac14(m-1)(M-1)+n+N\\
2\beta&=\tfrac14(m-1)(M-1)+n-N\\
2\gamma&=\tfrac14(m-1)(M-1)-n+N\\
2\delta&=\tfrac14(m-1)(M-1)-n-N
\end{aligned}
\]

% A544-P0064-C12
quarum quaelibet theorematis veritatem monstrat.

% A544-P0064-C13
\Article{3}

% A544-P0064-C14
Quodsi iam supponimus, $m$ et $M$ esse numeros primos, e combinatione theorematis praecedentis cum lemmate art. 1 theorema fundamentale protinus demanabit.\quad Patet enim,

% A544-P0064-C15
I.\quad quoties uterque $m,M$, sive alteruter tantum, sit formae $4k+1$, numerum $\frac14(m-1)(M-1)$ fore parem, adeoque $n$ et $N$ vel simul pares vel simul impares, et proin vel utrumque $m$ et $M$ alterius residuum quadraticum, vel utrumque alterius non-residuum quadraticum.

% A544-P0064-C16
II.\quad Quoties autem uterque $m,M$ est formae $4k+3$, erit $\frac14(m-1)(M-1)$ impar, hinc unus numerorum $n,N$ par, alter impar, et proin unus numerorum $m,M$ alterius residuum quadraticum, alter alterius non-residuum quadraticum. Q. E. D.

